\paragraph{Aerodynamics and the Betz limit.} Fig.~\ref{fig:A1_cp_curves} shows the Heier $C_p(\lambda,\beta)$ curves for $\beta=0$ to $20^\circ$. At $\beta=0$ the coefficient peaks at $C_{p,max}=0.480$ at $\lambda_{opt}=8.10$, which is 81.0 percent of the Betz limit $16/27=0.593$. The gap is the part of the wind power that real blades lose to drag, tip losses and wake rotation, none of which the actuator-disc argument contains. The turbine reaches its rated 2 MW at 12.00 m/s and 26.22 rpm (2.746 rad/s); the optimal-torque constant is $K_{opt}=96\,616$ W/(rad/s)$^3$ and the rotor inertia is $J=1.592\times10^{6}$ kg\,m$^2$ ($H=3$ s). Below rated wind speed the converter realises the $C_p$-optimal line by controlling the electrical power as $P_e=K_{opt}\omega^3$ (Fig.~\ref{fig:A2_power_speed}, dashed line). This is electrical power control through the converter, and it keeps the rotor on top of the $C_p$ hill. A rotor parked there has nothing to give: releasing its kinetic energy slows it to the left of the peak, where $C_p$ and aerodynamic power fall, so the support must be repaid afterwards. This is the origin of the secondary dip met in Task B.

\paragraph{Deloading, pitch and the reserve.} For a 10 percent reserve the converter instead holds the rotor to the right of the peak, at $C_p(\lambda_{del})=0.9\,C_{p,max}=0.432$. At 9 m/s this means $\omega_{del}=23.28$ rpm against $\omega_{opt}=19.67$ rpm ($\lambda_{del}=9.59$ against 8.10) and $P_{del}=759.4$ kW against 843.8 kW available, a reserve of 84.4 kW. The worked example $P=\tfrac12\rho A v^3C_p$ with $\rho=1.225$ kg/m$^3$, $A=3937$ m$^2$ ($R=35.4$ m) reproduces the 0.844 MW. On the overspeed side of the peak (Fig.~\ref{fig:A2_power_speed}, green line) $dP_{aero}/d\omega<0$: if the converter draws more than the rotor receives, the rotor decelerates, $\lambda$ moves back toward $\lambda_{opt}$ and $C_p$ and $P_{aero}$ rise. This is the physical basis of the project: kinetic energy and reserve are released together. Overspeed is used from 6 to 11 m/s. From the rated wind speed of 12 m/s, where the available power is capped at 2 MW, the rotor is held at its rated speed and the reserve is created by pitch control instead: $\beta=1.18^\circ$, $4.16^\circ$ and $8.48^\circ$ at 12, 13 and 14 m/s, holding 1.8 MW and keeping a 200 kW reserve (Fig.~\ref{fig:A3_power_curve}). Pitch actuators are slow, so pitch carries the steady reserve and the converter provides the fast response.

\paragraph{Reserve, stored energy and cost.} Fig.~\ref{fig:A4_reserve} shows what the grid can draw. The reserve grows from 25.0 kW at 6 m/s to 154 kW at 11 m/s and 200 kW at rated wind speed. The kinetic energy that can be released down to the 0.5 pu speed limit grows from 0.60 MJ (equivalent inertia $H_{eq}=0.30$ s) at 6 m/s to 4.34 MJ ($H_{eq}=2.17$ s) at 10 m/s and 4.50 MJ ($H_{eq}=2.25$ s) at 12 m/s. At 9 m/s one turbine releases 1.357 MJ between the deloaded speed and the MPPT speed ($H_{eq}=0.68$ s) and 3.232 MJ down to the speed limit ($H_{eq}=1.62$ s). The price is energy. For an illustrative Indian onshore site (Weibull $k=2$, $c=7.5$ m/s, Fig.~\ref{fig:A5_economics}) a permanent 10 percent deload cuts the annual energy from 4,725 MWh to 4,253 MWh (capacity factor 27.0 to 24.3 percent), a loss of 473 MWh or INR 16.5 lakh per year at 3.5 INR/kWh. This is an upper bound, because the reserve need only be held in low-inertia hours.
